% Figure: three models of the same pendulum, arranged by abstraction
% level. The axes are not numerical; they are conceptual (parameter
% count vs domain of validity).
\begin{figure}[htbp]
\centering
\begin{tikzpicture}[
  modelbox/.style={draw=engblack, fill=englime!18, rounded corners=2pt,
    minimum width=4.0cm, minimum height=1.1cm, font=\small, align=center},
  axislab/.style={font=\small\itshape, color=engblack!70},
  arrow/.style={-{Stealth}, thick, draw=engblack!70},
]
% Axes (conceptual)
\draw[arrow] (0, 0) -- (10.0, 0) node[right, axislab] {parameter count};
\draw[arrow] (0, 0) -- (0, 5.0) node[above, axislab, align=center] {domain\\ of validity};

% Three models
\node[modelbox] (m1) at (1.8, 1.0) {Linear small-angle\\
  $\ddot\theta + (g/L)\theta = 0$\\ (2 parameters)};
\node[modelbox] (m2) at (5.0, 2.6) {Exact small-amplitude\\
  $T = 4\sqrt{L/g}\,K(\sin(\theta_0/2))$\\ (3 parameters)};
\node[modelbox] (m3) at (8.2, 4.1) {Damped driven\\
  $\ddot\theta + 2\zeta\omega_0\dot\theta + \omega_0^2\theta = F/(mL)$\\
  (5 parameters)};

% Linking arrows (one model added a feature)
\draw[arrow, dashed, draw=engred!70] (m1) -- node[midway, fill=white,
  inner sep=2pt, font=\scriptsize\itshape, align=center] {drop small-\\angle
  assumption} (m2);
\draw[arrow, dashed, draw=engred!70] (m2) -- node[midway, fill=white,
  inner sep=2pt, font=\scriptsize\itshape, align=center] {add damping\\
  and forcing} (m3);

% Annotation
\node[align=left, font=\footnotesize, color=engblack!70] at (5.0, -0.9)
  {Each step adds parameters and extends the model's domain.
   The right model is the simplest model that answers the question.};
\end{tikzpicture}
\caption[Three pendulum models, arranged by abstraction]{Three models
of the same physical pendulum, arranged by parameter count and domain
of validity. The linear small-angle model has two parameters and a
narrow domain (small amplitude, undriven, undamped); the exact
small-amplitude model relaxes the linearisation at the cost of a third
parameter (amplitude $\theta_0$); the damped driven model extends to
forced and damped regimes with two further parameters. No model in
the figure dominates the others; each is the correct abstraction for
a question whose precision and regime it serves. The discipline of
abstraction is to choose the smallest model on the chain that the
question requires.}
\label{fig:vol02:ch18:abstraction-hierarchy}
\end{figure}
